\documentclass[11pt,a4paper]{article}
\usepackage[T1]{fontenc}
\usepackage[utf8]{inputenc}
\usepackage{lmodern}
\usepackage[margin=2.3cm]{geometry}
\usepackage{amsmath,amssymb}
\usepackage{booktabs}
\usepackage{graphicx}
\usepackage{xcolor}
\usepackage[colorlinks=true,linkcolor=blue!50!black,urlcolor=blue!50!black]{hyperref}
\usepackage{caption}
\usepackage{float}
\captionsetup{font=small,labelfont=bf}

\newcommand{\pass}{\textcolor{green!45!black}{\textbf{PASS}}}
\newcommand{\fail}{\textcolor{red!70!black}{\textbf{FAIL}}}
\newcommand{\code}[1]{\texttt{\hyphenchar\font=`\-#1}}
\sloppy
\emergencystretch=2em

\title{Pair functions in a hadron resonance gas at rest:\\ results of milestone M9
(spherical-harmonic anisotropy)}
\author{THERMINATOR 2 (hboltz fork), \code{src/work/hrg\_pair\_density\_spec.md}, Sec.~10}
\date{29 October 2026}

\begin{document}
\maketitle

\begin{abstract}
\noindent
Milestone M9 extends the Legendre moments of M2/M3 to the full decay direction
$\Omega^*=(\theta^*,\phi^*)$ in the helicity frame, with the 24 real spherical harmonics
$Y_{\ell m}$, $1\le\ell\le4$, at the ancestry levels L0, L1, L2, all ordered pairs and both weak
settings. The MC events are those of M2 (same seeds), which T9.1 confirms bit for bit. All tests
T9.0--T9.4 pass for both weak settings, with one deviation from the spec: the pooled pull-width
criterion is applied to bins with at least 500 entries instead of 50. With 50 entries the pooled
T9.2 widths are $1.006$--$1.012$ at every level, including L0. This excess is the Student-$t$ bias of a
pull whose error is estimated from the bin itself; it is visible only because pooling gives
$\sim10^5$ bins and a width tolerance of 0.3--0.6\%. The physics output: at L2 the only non-zero
moments are $a_{20}$ and $a_{40}$. $a_{40}$ is not negligible: $|a_{40}/a_{20}|$ grows from
$\sim0.05$ to $0.3$ with $\kappa$ in the $\pi\pi$ continuum at 0.5--0.7~GeV, and near the $\pi\pi$
threshold $a_{40}$ is as large as $a_{20}$ or larger. Truncating the chain anisotropy at an
$\epsilon_Q^2$ term is therefore adequate only away from threshold and at moderate $\kappa$.
\end{abstract}

\section{Spherical harmonics and moments}

\paragraph{Definitions.} With $x=\cos\theta$ and the associated Legendre functions without the
Condon--Shortley phase, $P_\ell^m(x)=(1-x^2)^{m/2}\,d^mP_\ell/dx^m$ ($0\le m\le\ell$, $P_\ell^0=P_\ell$),
the real spherical harmonics are
\begin{equation}
Y_{\ell0}=N_{\ell0}P_\ell(\cos\theta),\qquad
\begin{Bmatrix}Y^{c}_{\ell m}\\ Y^{s}_{\ell m}\end{Bmatrix}=\sqrt2\,N_{\ell m}P_\ell^m(\cos\theta)
\begin{Bmatrix}\cos m\phi\\ \sin m\phi\end{Bmatrix},\qquad
N_{\ell m}=\sqrt{\frac{2\ell+1}{4\pi}\frac{(\ell-m)!}{(\ell+m)!}},
\end{equation}
orthonormal on the sphere, $\int d\Omega\,Y_{\ell m}Y_{\ell'm'}=\delta_{\ell\ell'}\delta_{mm'}$. For
$1\le\ell\le4$ there are $2\ell+1$ functions per $\ell$, 24 in total.

\paragraph{$m=0$: the Legendre moments.} $Y_{\ell0}=\sqrt{(2\ell+1)/4\pi}\,P_\ell(\cos\theta)$ does not
depend on $\phi$, so $a_{\ell0}=\sqrt{(2\ell+1)/4\pi}\,\langle P_\ell\rangle$: the $m=0$ moments are the
$a_\ell$ of M2/M3 in another normalization ($\int_{-1}^1P_\ell^2dx=2/(2\ell+1)$ against
$\int Y_{\ell0}^2d\Omega=1$). They describe a distribution that is axially symmetric around the polar
axis $\hat z^*=\hat{\mathbf K}$. In the gas at rest the fluid velocity in the pair frame points along
$-\hat{\mathbf K}$ and $\epsilon_Q=2\kappa p^*\cos\theta^*/M$ (equal masses) depends on $\theta^*$ only.
So every $m\neq0$ moment vanishes, which is the content of T9.2. An $\epsilon_Q^n$ dependence of the pair
functions contributes only $\ell\le n$ with the parity of $n$. An $\epsilon_Q^2$ term,
$\epsilon_Q^2=c^2[1+2P_2(\cos\theta^*)]/3$ with $c=2\kappa p^*/M$, contributes to the $\Omega^*$-averaged density
($\ell=0$, the $c^2/3$ part, contained in the L2 table; $a_{00}=1/\sqrt{4\pi}$ itself is fixed by the
normalization) and to $a_{20}$. $a_{40}\neq0$ needs $\epsilon_Q^4$ or higher.

\paragraph{$m\neq0$: flow.} In a moving gas the pair frame (pure boost along $\mathbf K$) keeps the
part of the fluid four-velocity transverse to $\mathbf K$, $\mathbf u^*_\perp=\mathbf u_{\perp}^{\rm lab}$, so the
axis $\hat u^*$ entering $\epsilon_Q=-2p^*\,\hat n\cdot\mathbf u^*$ is tilted away from $\hat{\mathbf K}$.
In one fluid cell the distribution is then axially symmetric around $\hat u^*$,
$\sum_\ell c_\ell P_\ell(\hat n\cdot\hat u^*)$, and the addition theorem
$P_\ell(\hat n\cdot\hat u^*)=\frac{4\pi}{2\ell+1}\sum_mY_{\ell m}(\hat n)Y_{\ell m}(\hat u^*)$
moves strength from $a_{\ell0}$ into $a_{\ell m}$, $m\neq0$, of the \emph{same} $\ell$, with a $\phi^*$
dependence set by the azimuth of $\hat u^*$ relative to $\hat x^*$ (the beam). Averaged over cells,
reflection symmetries remove the $\sin m\phi^*$ terms for a central source and the odd-$m$ $\cos$ terms at
midrapidity; $\cos2\phi^*$, which measures $\langle u^{*2}_x\rangle-\langle u^{*2}_y\rangle$, generally
survives. L0 and L1 stay isotropic with any flow (one-step decays with constant matrix elements are
isotropic in the pair rest frame), so their $a_{\ell m}$ remain null tests.

\paragraph{Estimating the coefficients.} The pair distribution in a bin $(M,\kappa)$ is expanded as
\begin{equation}
\frac{1}{N}\frac{dN}{d\Omega^*}=\frac{1}{4\pi}+\sum_{\ell\ge1}\sum_m a_{\ell m}\,Y_{\ell m}(\Omega^*),\qquad
a_{\ell m}=\int d\Omega^*\,Y_{\ell m}(\Omega^*)\,\frac1N\frac{dN}{d\Omega^*}=\langle Y_{\ell m}\rangle,
\end{equation}
by orthonormality. The MC estimator is the sample mean over the $N$ ordered pairs of the bin,
\begin{equation}
\hat a_{\ell m}=\frac1N\sum_{i=1}^NY_{\ell m}(\theta^*_i,\phi^*_i),\qquad
\sigma^2(\hat a_{\ell m})=\frac1N\Big(\frac1N\sum_iY_{\ell m}^2(\theta^*_i,\phi^*_i)-\hat a_{\ell m}^2\Big),
\end{equation}
which is what a \code{TProfile2D} filled with $Y_{\ell m}$ stores (sums of $Y$ and $Y^2$ and the entries
$N$; the error is the profile's default). For an isotropic distribution $\langle Y_{\ell m}^2\rangle=1/4\pi$,
so $\sigma\simeq1/\sqrt{4\pi N}$. Like-sign pairs enter twice with identical even-$\ell$ values, so their
errors are multiplied by $\sqrt2$. Cumulative levels and wider $(M,\kappa)$ cells add the sums before
dividing: $\hat a=\sum_bS_b/\sum_bN_b$ with $S_b=N_b\hat a_b$, and the same for $\sum Y^2$. The angles of
pair $i$ are $\cos\theta^*_i=\hat n_i\cdot\hat z^*$ and $\phi^*_i=\mathrm{atan2}(\hat n_i\cdot\hat y^*,\hat n_i\cdot\hat x^*)$
in the frame of Sec.~2.

\section{Implementation}

\paragraph{Harmonics.} \code{RealYlm(cth, phi, y)} in \code{PairFunctions.h}: $P_\ell^m$ by the
standard recursion in $\ell$ at fixed $m$ starting from $P_m^m=(2m-1)!!\sin^m\theta$ (no
Condon--Shortley phase), $N_{\ell m}$ of spec Sec.~10.2, $\sqrt2$ for $m\neq0$. Index
$\ell^2-1+k$ with $k=0$ ($m=0$), $2m-1$ ($\cos m\phi$), $2m$ ($\sin m\phi$); tags \code{10},
\code{21c}, \code{21s}, \dots

\paragraph{MC.} \code{ex\_pair\_functions\_mc.exe --harmonics} fills, per exclusive level (L0, L1x,
L2x) and ordered pair, 24 \code{TProfile2D} \code{hY<tag>\_<level>\_<pair>} with the $(M,\kappa)$
binning of \code{hP1}, \code{hP2}. Frame (spec Sec.~10.1): $\hat z^*=\hat{\mathbf K}$,
$\hat x^*$ = lab $\hat z$ made transverse to $\hat{\mathbf K}$ (lab $\hat x$ if the norm is below
$10^{-6}$), $\hat y^*=\hat z^*\times\hat x^*$, $\hat n$ = direction of $\mathbf p_1^*$ after the
pure boost; $\cos\theta^*$ is computed exactly as for \code{hP1}. The option draws no random numbers.
Runs: $8\times2.5\times10^4$ events per weak setting with the M2 seeds (1--8 weak on, 101--108 weak
off) in \code{build/t2/}, with the same \code{Integrator} cache; about 35~ms per event (twice the
M2 cost) and 0.5~GB per job.

\paragraph{Analysis.} \code{ex\_pair\_functions\_harmonics.exe} runs T9.0 alone, or T9.0--T9.4 and the
report tables on a merged file (\code{--mc}, \code{--ref} for T9.1, \code{--minentries}). Pull
criterion as T2.3: $|\text{mean}|<3/\sqrt n$ and $|\text{width}-1|<3/\sqrt{2n}$ over bins with at
least \code{--minentries} entries. The errors of like-sign even-$\ell$ moments are scaled by $\sqrt2$
(doublets). Each test family is run per (pair, cumulative level, function) and pooled per (level,
function) and per level. Pooling leaves out the mirror pairs $\pi^-\pi^+$, $K^-K^+$, $\bar pp$,
which are the $(+-)$ pairs with $\hat n\to-\hat n$ (identical even-$\ell$ and opposite odd-$\ell$
moments). These pairs are still tested individually. With several hundred individual tests at a
false-alarm rate of $\approx0.5\%$, a few chance failures are expected. A family passes if every
pooled test passes and the number of failed individual tests is consistent with the false-alarm
rate (Poisson $p>10^{-3}$). The rate is computed exactly for Gaussian pulls: the mean and the
width are independent, with $n\,\text{width}^2\sim\chi^2_{n-1}$.

\section{Results}

\begin{table}[H]
\centering
\small
\begin{tabular}{llll}
\toprule
Test & Content & weak on & weak off \\
\midrule
T9.0 & orthonormality (32$\times$64 quadrature) & $5.8\times10^{-15}$ \pass & same \\
     & parity $Y(-\hat n)=(-1)^\ell Y(\hat n)$, $10^4$ directions & $2.5\times10^{-15}$ \pass & same \\
T9.1 & entries vs M2 profiles & identical, 241\,012 bins & identical, 237\,326 bins \\
     & $a_{10},a_{20}$ vs $\sqrt{(2\ell+1)/4\pi}\langle P_\ell\rangle$ & $6.4\times10^{-14}$ \pass &
       $5.8\times10^{-14}$ \pass \\
\midrule
\multicolumn{4}{l}{\emph{bins with $\ge500$ entries (verdict)}} \\
T9.2 & $m\neq0$; pooled width L0 / L1 / L2 & 1.000 / 1.002 / 1.000 & 1.002 / 1.003 / 1.001 \\
     & individual failures (expected) & 0 / 264 (1.4) \pass & 0 / 264 (1.4) \pass \\
T9.3 & odd $\ell$, $m=0$, unlike-sign; pooled width & 1.001 / 1.000 / 1.001 & 1.000 / 0.995 / 0.996 \\
     & individual failures & 0 / 24 \pass & 0 / 24 \pass \\
T9.4 & $\ell=2,4$, $m=0$, L0 / L1; pooled width & 1.000 / 1.000 & 0.990 / 0.989 \\
     & individual failures & 0 / 16 \pass & 0 / 16 \pass \\
\midrule
\multicolumn{4}{l}{\emph{bins with $\ge50$ entries (spec)}} \\
T9.2 & pooled width L0 / L1 / L2 & 1.006 / 1.007 / 1.006 & 1.009 / 1.011 / 1.009 \\
     & individual failures (expected) & 3 / 336 (1.8) & 6 / 336 (1.8) \\
     & verdict & \fail{} (pooled widths) & \fail{} (pooled widths) \\
T9.3, T9.4 & & \pass & \pass \\
\bottomrule
\end{tabular}
\caption{M9 tests. Pooled means per level at most 0.004 (T9.2), 0.012 (T9.3), 0.016 (T9.4) in
magnitude; per (level, function) at most 0.043, all within their tolerances. With 500 entries the like-sign pairs drop out of T9.4 at L0
and L1 (no bins left), hence 16 instead of 24 tests. Logs: \code{m9\_weak*.log} (50 entries),
\code{m9\_weak*\_min500.log} (500).}
\end{table}

\paragraph{T9.1.} The new runs reproduce the M2 events: entries are identical in every bin of the
$a_{10}$ and $a_{20}$ profiles at all three exclusive levels, and the moments agree to $6\times10^{-14}$
relative to $\max(|a|,\sigma)$. A pure relative difference would be meaningless in bins with mean
$\approx0$. The pair counts of the 16 job logs also equal those of the M2 jobs seed by seed.

\paragraph{Pull widths at 50 entries.} Table~\ref{tab:width} splits the T9.2 pulls by the
number of entries of the bin. The excess is confined to bins with fewer than $\sim200$ entries and
follows the Student-$t$ expectation $\sqrt{N/(N-3)}$ for a pull whose error is the profile's
spread estimate, $s/\sqrt N$ with the $1/N$ variance. For $N\ge1000$ the width is 1 within its
error at every level, L2 included. So there is no sign of correlated entries (pairs sharing a
particle in a cascade) beyond the like-sign doublets. M3 (T2.3, T3.1) did not see this because it
pooled far fewer bins (tolerance 2--5\%).

\begin{table}[H]
\centering
\small
\begin{tabular}{lrrrr}
\toprule
entries & L0 width & L2 width & $\pm$ & $t$-expectation \\
\midrule
50--100 & 1.020 & 1.027 & 0.004 / 0.003 & 1.025 / 1.031 \\
100--200 & 1.002 & 1.014 & 0.006 / 0.003 & 1.012 / 1.018 \\
200--500 & 1.006 & 1.004 & 0.007 / 0.003 & 1.005 / 1.007 \\
$10^3$--$10^4$ & 0.997 & 0.999 & 0.004 / 0.002 & 1.001 / 1.001 \\
$10^4$--$10^5$ & 1.001 & 1.001 & 0.004 / 0.002 & 1.000 / 1.000 \\
\bottomrule
\end{tabular}
\caption{T9.2 pull width versus bin entries, weak on (\code{width\_vs\_entries.C},
\code{width\_vs\_entries\_weak*.txt}; weak off the same). The $t$-expectation is averaged over the
bins of the class (like-sign: $N/2$ independent entries). The measured excess is slightly below it
at the lowest $N$ because $Y_{\ell m}$ is bounded and has lighter tails than a Gaussian.}
\label{tab:width}
\end{table}

\begin{figure}[H]
\centering
\includegraphics[width=\textwidth]{fig/pulls.png}
\caption{Pull distributions of T9.2--T9.4 (bins with $\ge500$ entries, mirror pairs excluded), with
the unit Gaussian.}
\end{figure}

\section{Report: $a_{20}$ and $a_{40}$ at L2}

The only moments that are non-zero at L2 are the even-$\ell$, $m=0$ ones, as expected.
Figures~\ref{fig:pipm}--\ref{fig:kk} show $a_{20}$ and $a_{40}$ at L2 versus $\kappa$ at the M3 test
masses, with $a_{20}$ at L1 for contrast. The $(M,\epsilon_K)$ cell tables (as in M3) are the
\code{M9tab} lines of the logs. Highlights, weak on (weak off similar):

\begin{itemize}
\item $\pi^+\pi^-$, $M=0.5$--$0.7$~GeV: $a_{20}$ from $-0.004$ ($\kappa<0.2$) to $-0.083$
  ($\kappa=1$--2~GeV) and $a_{40}$ from 0 to $+0.023$, so $a_{40}/a_{20}$ goes from $-0.04$ to
  $-0.27$ (weak off $-0.06$ to $-0.28$). Consistent with M3: $a_{20}=\sqrt{5/4\pi}\,a_2$.
\item $\pi^+\pi^-$ near threshold, $M=0.28$--$0.40$~GeV: $a_{20}$ changes sign inside the band
  (at $M=0.30$ it rises to $+0.06$ at large $\kappa$, while the cell average is $-0.004$ to $-0.009$).
  $a_{40}$ reaches $-0.03$, as large as $a_{20}$ or larger: $a_{40}/a_{20}$ ranges from $+0.9$ to $-0.5$
  along $\kappa$ at $M=0.30$.
\item $\pi^+\pi^-$, $0.40$--$0.50$ and $0.70$--$0.85$~GeV: $|a_{40}/a_{20}|\lesssim0.1$.
\item $\pi^+\pi^+$: same pattern, larger: at 0.5--0.7~GeV $a_{20}$ down to $-0.10$ ($-0.15$ weak off),
  $a_{40}/a_{20}$ down to $-0.34$ ($-0.37$); at threshold $a_{40}$ dominates (down to $-0.068$
  at $\kappa>1$~GeV).
\item $K^+K^-$: $\phi$ pole consistent with 0; continuum 1.03--1.20~GeV $a_{20}$ up to $+0.033$,
  $a_{40}/a_{20}\approx-0.05$ to $-0.18$; above 1.2~GeV consistent with 0 within $\sim0.002$.
\item Pole bins (one 4-MeV bin each) contain the pole and the L2 continuum underneath, so their
  moments are those of the continuum diluted by the pole fraction. Where the pole dominates they
  are consistent with 0: $\rho^0$, $f_0$, $\phi$ ($|a_{20}|,|a_{40}|\lesssim2\times10^{-3}$, the
  largest values at the edges of the $\kappa$ range). For the $\omega$ (BR$(\pi^+\pi^-)=1.5\%$),
  $a_{20}=-0.008$ to $-0.017$ at $\kappa>0.4$~GeV (10--15$\sigma$). For the $K^0_S$ with weak decays on,
  $a_{20}\approx-0.001$ ($\sim5\sigma$); with weak off the $K^0$ does not decay, the bin is pure
  continuum and $a_{20}$ reaches $-0.08$. Like-sign pairs have no pole: there the "pole bins" are
  continuum ($a_{20}$ down to $-0.13$).
\end{itemize}

\noindent So an $\epsilon_Q^2$ (pure $\ell=2$) term describes the chain anisotropy to about 10\% only
for $\kappa\lesssim0.7$~GeV away from threshold. At large $\kappa$ the $\ell=4$ component is 20--40\%
of the $\ell=2$ one, and near the $\pi\pi$ threshold it is the leading term.

\begin{figure}[H]
\centering
\includegraphics[width=\textwidth]{fig/moments_pippim_weakon.png}
\caption{$\pi^+\pi^-$, weak on: $a_{20}$, $a_{40}$ at L2 and $a_{20}$ at L1 versus $\kappa$, at the
$\rho^0$ pole bin and in 20-MeV continuum windows (5 $M$ bins) at 0.30, 0.45, 0.60~GeV.}
\label{fig:pipm}
\end{figure}

\begin{figure}[H]
\centering
\includegraphics[width=0.78\textwidth]{fig/moments_pippip_weakon.png}
\caption{$\pi^+\pi^+$, weak on (errors include the doublet factor $\sqrt2$).}
\end{figure}

\begin{figure}[H]
\centering
\includegraphics[width=0.78\textwidth]{fig/moments_KpKm_weakon.png}
\caption{$K^+K^-$, weak on: $\phi$ pole bin and continuum windows at 1.10, 1.30~GeV.}
\label{fig:kk}
\end{figure}

\noindent Weak-off figures: \code{fig/moments\_*\_weakoff.png}.

\section{Outputs and status}

\begin{itemize}
\item MC: \code{build/t2/m9/harm\_weak\{on,off\}.root} (hadd of \code{m9/jobs/}), driver
  \code{m9/run\_jobs.sh}.
\item Analysis: \code{m9/harm\_weak*\_ana.root} (50 entries), \code{\_ana500.root} (500):
  \code{pulls\_T9\{2,3,4\}}, \code{a20\_<lev>\_<pair>\_<mass>}, \code{a40\_\dots},
  \code{a20\_cells\_<lev>\_<pair>}, \code{a40\_cells\_\dots}.
\item Deviation from the spec: the pull-width verdict uses bins with $\ge500$ entries; the
  50-entry run is reported alongside.
\end{itemize}

\end{document}
